\section*{Planteamiento del Proyecto}
\justifying

El proyecto consiste en el diseño y construcción de un sistema didáctico modular con el objetivo de
facilitar la enseñanza y el aprendizaje de la electrónica de potencia. Este sistema permitirá a 
estudiantes o entusiastas experimentar con 
diversos módulos, tales como rectificadores, inversores, convertidores DC-DC, entre otros.

Mediante la observación de señales de entrada y salida y el ajuste de parámetros en tiempo real,
el sistema permitirá a los usuarios aprender y experimentar con diversos conceptos de
electrónica de potencia, promoviendo así un aprendizaje activo e interactivo.

Además, contará con un manual de usuario que explique el funcionamiento teórico de cada módulo y
una guía práctica que oriente sobre los parámetros a modificar y los efectos a analizar. Con 
este enfoque, se espera que los usuarios desarrollen una comprensión profunda y tangible de los 
principios de la materia.

El diseño modular del sistema permitirá que, en el futuro, se puedan incorporar
expansiones o módulos personalizados según necesidades del usuario o avances tecnológicos.
De este modo, el proyecto queda abierto a mejoras continuas y a la adaptación a nuevas áreas
de estudio de la electrónica.

El sistema propuesto surge de la necesidad de contar con herramientas para el aprendizaje de la
electrónica de potencia que sean accesibles, seguras y efectivas. Desarrollar un sistema didáctico
modular permitirá la enseñanza y el aprendizaje de los conceptos de la cátedra, sin necesidad de
contar con un laboratorio especializado y costoso. Asimismo, se podrán visualizar de manera clara
y simple los efectos de cambiar parámetros en los distintos circuitos, agilizando los tiempos de
aprendizaje, en especial debido al poco tiempo disponible en el año que se tiene para la enseñanza
de la materia, teniendo en cuenta el amplio contenido que se debe abordar (principal problema
en la cátedra según el titular Ing. Alejandro Giorgis) 

Algunos desafíos que se presentan en el desarrollo de este proyecto son:

\begin{enumerate}[noitemsep]
    \item Selección de módulos a implementar.
    \item Selección de componentes y materiales.
    \item Diseño de circuitos y placas, con cálculos de potencia y disipación de calor.
    \item Seguridad para garantizar la integridad de los usuarios y del sistema.
    \item Asegurar el fácil mantenimiento y reparación del sistema, en caso de ser necesario.
    \item Diseño de una interfaz de usuario intuitiva y accesible.
    \item Desarrollo de un manual de usuario claro y detallado.
    \item Adaptabilidad para futuras expansiones y mejoras.
    \item Diseño de los módulos para que sean fáciles de intercambiar y conectar.
\end{enumerate}